\documentclass[10pt]{article}
\newcommand{\documenttype}{Document}
\usepackage[a4paper,margin=0.73in]{geometry}
\usepackage[T1]{fontenc}
\usepackage[utf8]{inputenc}
\usepackage{lmodern,amsmath,amssymb,microtype,xcolor,enumitem,tabularx,titlesec,fancyhdr,hyperref}
\definecolor{accent}{HTML}{315B8A}
\definecolor{link}{HTML}{5B4F7A}
\hypersetup{colorlinks=true,urlcolor=link,linkcolor=link,citecolor=link,pdfauthor={Phurailatpam Hemantakumar}}
\setlength{\parindent}{0pt}\setlength{\parskip}{5pt}\setlength{\headheight}{14pt}\setlength{\emergencystretch}{2em}
\titleformat{\section}{\normalsize\bfseries\color{accent}}{}{0pt}{}
\titlespacing*{\section}{0pt}{6pt}{3pt}
\setlist[itemize]{leftmargin=16pt,itemsep=2pt,topsep=3pt,parsep=0pt}
\pagestyle{fancy}\fancyhf{}\fancyhead[L]{\small Phurailatpam Hemantakumar}\fancyhead[R]{\small\documenttype}\fancyfoot[R]{\small\thepage}\renewcommand{\headrulewidth}{0.3pt}
\newcommand{\firstheading}[2]{\thispagestyle{plain}{\LARGE\bfseries\color{accent}#1}\par\vspace{3pt}{#2}\par\vspace{4pt}{\color{accent}\hrule}\vspace{8pt}}
\newcommand{\dated}[2]{\par\noindent\begin{tabularx}{\linewidth}{@{}X r@{}}#1&#2\end{tabularx}\par}
\newcommand{\contact}{\href{mailto:hemantaph@link.cuhk.edu.hk}{hemantaph@link.cuhk.edu.hk}\quad|\quad +852 9867 8178\quad|\quad\href{https://hemantaph.com}{hemantaph.com}\\Department of Physics, The Chinese University of Hong Kong}
\renewcommand{\documenttype}{Selection criteria}
\hypersetup{pdftitle={Response to key selection criteria | Monash OzGrav 698908}}
\setlength{\parskip}{5pt}
\begin{document}
\firstheading{Response to key selection criteria}{Phurailatpam Hemantakumar\quad|\quad Research Fellow -- Gravitational Wave Discovery (OzGrav)\\Monash University, Job No. 698908\quad|\quad 21 November 2026}

\section*{1. A doctoral qualification in the relevant discipline or a closely related field}
I have completed the research and examination requirements for a PhD in Physics at The Chinese University of Hong Kong (CUHK), supervised by Dr Otto Akseli Hannuksela. My thesis, \textit{Statistical and Data Analysis of Eccentric and Lensed Gravitational Waves}, investigates statistical predictions for gravitational-wave (GW) lensing and the modelling and inference of eccentric compact-binary signals. The formal degree award is expected in December 2026. My earlier qualifications are an M.Sc. in Physics from IIT Bombay and a B.Sc. (Hons.) in Physics from the University of Delhi.

\section*{2. Experience in gravitational-wave astronomy, general relativity, or high-energy astrophysics (or a closely related field)}
My research is in GW astronomy, particularly strong gravitational lensing, compact-binary populations, signal detectability, and waveform modelling. I lead development of \href{https://ler.hemantaph.com/}{\texttt{ler}}, which models lensed and unlensed compact-binary populations and propagates source, lens, and detector selection into detectable-event forecasts. I also lead development of \href{https://gwsnr.hemantaph.com/}{\texttt{gwsnr}} for efficient signal-to-noise ratio and detection-probability calculations. As a member of the LIGO Scientific Collaboration, I contribute lensing statistics to analyses of observed events and LVK population studies.

I co-developed \texttt{TaylorF2Ecck}, an analytic post-Newtonian model for low-eccentricity inspirals, implemented it in a \texttt{LALSuite} fork, and performed Bayesian parameter estimation for GW170817 and GW190425. Related work investigated the eccentricity--line-of-sight-acceleration degeneracy for GW190814. This work combines general-relativistic signal modelling, detector time-series analysis, and Bayesian inference.

\section*{3. Demonstrated analytical and manuscript preparation skills, including a track record of refereed research publications}
My analytical work connects physical models to testable statistical predictions. For \texttt{TaylorF2Ecck}, I contributed waveform implementation, consistency checks, injections, and parameter-estimation analyses used to assess eccentricity constraints in inspiral-dominated signals. For GW lensing, I developed Monte Carlo implementations of elliptical power-law lenses with external shear and numerical cross-sections, and studied the impact of population and detector selection on detection rates and lensing priors.

I am lead author of refereed papers in \textit{Physical Review D} (2026; eccentric waveforms) and the \textit{Journal of Open Source Software} (2026; \texttt{ler}). My refereed co-authored work includes articles in \textit{The Astrophysical Journal}, \textit{Monthly Notices of the Royal Astronomical Society}, and \textit{Philosophical Transactions of the Royal Society A}. I have contributed to manuscript development, technical validation, responses to collaborators, and public scientific documentation. A selected publication list is included in my CV, with a full list available separately.

\section*{4. Ability to solve complex problems using discretion, innovation and diagnostic skills and/or expertise}
The main computational challenges in my research arise from combining physical accuracy, statistical selection, and numerical efficiency. In \texttt{ler}, I developed calculations for complex lens potentials and their image properties, including optical depths and numerical strong-lensing cross-sections. I use targeted numerical checks and comparisons of independently implemented methods to diagnose inconsistencies. In \texttt{gwsnr}, I implemented partial-scaling interpolation and accelerated waveform inner products using Numba, JAX, and MLX to reduce the cost of large population simulations while retaining a direct-calculation baseline for validation.

In waveform analysis, I use injection studies and comparisons with alternate approximants to identify where low-eccentricity expansions become inaccurate. These examples illustrate my approach to difficult problems: identify the physical and numerical approximations, design diagnostic tests, quantify their limitations, and revise the implementation accordingly.

\newpage
\section*{5. Well-developed planning and organisational skills, including prioritising tasks and meeting deadlines}
I coordinate software development, validation, documentation, and scientific analyses across several collaborations. I have managed work on \texttt{ler} and \texttt{gwsnr} alongside my PhD, refereed publications, gravitational-wave analyses, student mentoring, and invited workshops. My ongoing LVK/GRASP responsibilities include planned \texttt{ler} and \texttt{gwsnr} contributions towards GWTC-6 and GWTC-7.

For research projects, I divide work into reproducible stages: method development, unit and injection tests, population-scale analysis, manuscript preparation, and release or documentation. This enables collaborators to assess intermediate results and helps me prioritise essential validations before extending a calculation. Completing two lead-author papers in 2026 while maintaining research software demonstrates my ability to deliver parallel commitments.

\section*{6. Excellent written and verbal communication, including clear and succinct reports and documents}
I have written and co-written refereed articles, software papers, scientific documentation, and research presentations for the GW community. I maintain user-facing documentation for \texttt{ler} and \texttt{gwsnr}, translating specialised numerical methods and modelling assumptions into examples that other researchers can reproduce. My talks include an East-Asia GW Science Workshop presentation on \texttt{ler} (Nagoya, 2025), a seminar on detectable lensed GW populations (IUCAA, 2026), and an LVK Collaboration Meeting poster on \texttt{TaylorF2Ecck} (Pune, 2026).

I have taught scientific Python at a CUHK summer school and conducted GW data-analysis workshops in India. These activities require adapting technical explanations to participants with different levels of experience, and responding constructively to questions and feedback.

\section*{7. Awareness of confidentiality, privacy and appropriate information handling}
Work within the LIGO Scientific Collaboration and multi-institutional research projects requires a clear distinction between public results and internal or unpublished analyses. I understand the importance of respecting collaboration review and release procedures, appropriately restricting circulation of non-public material, and giving due credit to colleagues. My software work similarly requires attention to code provenance, reproducible research outputs, and responsible management of shared research materials.

In a Monash appointment, I would follow applicable University and OzGrav requirements for information security, confidentiality, privacy, data access, and publication approval, seeking advice when the status of information is uncertain. I would not disclose unpublished collaborator results or restricted materials without the appropriate permissions.

\section*{8. Demonstrated capacity to work collegially with other staff}
My research is conducted across international teams. I lead \texttt{ler} development with colleagues at CUHK and other LVK institutions; I work with Achamveedu Gopakumar, Maria Haney, Tjonnie Li, and Srishti Tiwari on eccentric-waveform modelling; and I collaborate with Anupreeta More on multimessenger lensing. I am also working with researchers at Ewha Womans University on star-formation-informed compact-binary merger histories.

These projects require agreed interfaces between theory, software, simulation, and data analysis, as well as constructive discussion of different technical approaches. I have supervised or mentored undergraduate, master's, and early-stage PhD researchers, including co-authored outcomes. I would bring this collaborative approach to Monash's OzGrav node, contributing to shared analyses and software while developing a research programme complementary to colleagues working on GW populations, inference, and gravity.
\end{document}
